\documentclass[11pt]{article}

\usepackage[utf8]{inputenc}
\usepackage[T1]{fontenc}
\usepackage{lmodern}
\usepackage{geometry}
\usepackage{amsmath,amssymb,amsfonts,amsthm,mathtools}
\usepackage{bm}
\usepackage{enumitem}
\usepackage{microtype}
\usepackage{hyperref}
\usepackage{titlesec}
\usepackage{booktabs}
\usepackage{array}
\usepackage{setspace}
\usepackage{mathrsfs}
\usepackage{physics}

\geometry{margin=1in}
\hypersetup{colorlinks=true, linkcolor=black, urlcolor=blue, citecolor=black}
\setlist[enumerate]{topsep=0.4em,itemsep=0.35em}
\setlist[itemize]{topsep=0.3em,itemsep=0.25em}
\titleformat{\section}{\large\bfseries}{\thesection.}{0.5em}{}
\titleformat{\subsection}{\normalsize\bfseries}{\thesubsection.}{0.5em}{}
\setstretch{1.06}

\newcommand{\Aether}{\AE{}ther}
\newcommand{\Aflow}{\AE{}ther-flow}
\newcommand{\TheoryName}{\Aflow{} Interpretation of Relativity}
\newcommand{\TheoryProgram}{\Aether{} / \Aflow{} framework}
\newcommand{\Sub}{\mathcal{A}}
\newcommand{\BridgeMetric}{\mathfrak{g}}
\newcommand{\Ell}{\mathcal{L}}

\newtheorem{definition}{Definition}
\newtheorem{proposition}{Proposition}
\newtheorem{remark}{Remark}
\newtheorem{corollary}{Corollary}
\newtheorem{theorem}{Theorem}

\title{\textbf{The \TheoryName{}}\\[0.4em]
\large Higher-Derivative Quartic Branch Failure of the Unit-Lapse Coulomb-Split Regular-Part Closure Test}
\author{}
\date{}

\begin{document}

\maketitle

\begin{abstract}
The propagated trace-monopole cancellation test has now been carried out for the fixed-completion unit-lapse spatial-harmonic formulation, and at the present recorded scope it does not close: a vanishing initial trace monopole is not derived from the existing compatibility conditions, and its propagation is not presently proved. The next smallest repair is therefore the longitudinal auxiliary renormalization
\[
\chi=\chi_{\mathrm{Coul}}+\chi_{\mathrm{reg}},
\]
with the hope that the regular part might lie in the same unweighted Sobolev class used by the old local and coercive architecture.

This note performs that next test. The Coulomb split does remove the explicit monopole tail at the slice equation level: after choosing a compactly supported unit-mass profile \(B_{\mathrm{ext}}\), the regular part satisfies
\[
\kappa_\theta \Delta \chi_{\mathrm{reg}}
=
-\bigl(F_\chi-\mathfrak m_\chi B_{\mathrm{ext}}\bigr),
\qquad
\int_{\mathbb R^3}\bigl(F_\chi-\mathfrak m_\chi B_{\mathrm{ext}}\bigr)\,dx=0.
\]
Thus the leading \(1/|x|\) coefficient is subtracted explicitly.

The decisive point is that this is still not enough to recover the \emph{old} unweighted \(H^N\) closure package. For the zero-mean source
\[
F_{\chi,\mathrm{reg}}
\equiv
F_\chi-\mathfrak m_\chi B_{\mathrm{ext}},
\]
the Fourier representation of the regular part is
\[
\widehat{\chi_{\mathrm{reg}}}(\xi)
=
-\frac{\widehat{F_{\chi,\mathrm{reg}}}(\xi)}{\kappa_\theta |\xi|^2}.
\]
Hence unweighted \(H^N\) control of \(\chi_{\mathrm{reg}}\) requires an additional low-frequency seminorm
\[
\int_{|\xi|\le 1} |\xi|^{-4}\,|\widehat{F_{\chi,\mathrm{reg}}}(\xi)|^2\,d\xi,
\]
not merely zero mean and not merely the old \(H^{N-2}\) control of the source. A zero-mean condition removes only \(\widehat{F_{\chi,\mathrm{reg}}}(0)\); it does not control the remaining \(|\xi|^{-2}\) singularity of the Poisson inverse in unweighted \(L^2\). Therefore the Coulomb split does not presently close the old unweighted \((K,\chi,\beta)\) local/coercive architecture. The next honest move is now the broader weighted or homogeneous function-space reformulation whose controlled norm records precisely that missing low-frequency information.
\end{abstract}

\section{Introduction}

The fixed-completion unit-lapse spatial-harmonic route has now passed through two increasingly narrow repair tests. First, the straightforward unweighted \((K,\chi,\beta)\) package failed because the longitudinal auxiliary variable \(\chi\) acquires a trace-driven Coulomb tail whenever the source monopole
\[
\mathfrak m_\chi(t)
=
\int_{\mathbb R^3}\bigl(K-\mathcal N_\chi\bigr)(t,x)\,dx
\]
is nonzero. Second, the propagated trace-monopole cancellation test asked whether that monopole could be forced at \(t=0\) and then propagated by the existing unit-lapse evolution. At the present recorded scope, that route also failed to close.

The next smallest repair is therefore not yet a weighted-space redesign. It is the Coulomb split itself: subtract the explicit monopole profile from the longitudinal auxiliary variable and ask whether the remainder \(\chi_{\mathrm{reg}}\) fits back into the old unweighted Sobolev energy.

\paragraph{Framework and claim boundary.}
\TheoryProgram{} denotes the broader \Aether{} / \Aflow{} framework built on the claims that the \Aether{} is the underlying four-dimensional substrate of reality, the \Aflow{} is its intrinsic ordered motion, observed three-dimensional space is the local experiential slice of that deeper substrate, S-time is the experienced order of change arising from matter, light, and the \Aflow{}, and observed expansion is the three-dimensional appearance of deeper four-dimensional ordered motion. Within that framework, \TheoryName{} denotes the exact-closure theory in which the effective gravitational dynamics are adopted to be exactly Einsteinian with universal matter coupling. In the active sequence, \emph{adoption} means use of that established relativistic dynamics without claiming substrate derivation, while \emph{derivation} is reserved for a first-principles recovery from explicit substrate variables.

The present note stays within that boundary. It does not claim that no Coulomb-split theorem can ever be proved in any refined function class. It records only that the split does not presently restore the \emph{old unweighted} closure package.

\section{The Coulomb Split to Be Tested}

On each slice of the unit-lapse reduction, write the longitudinal source as
\begin{equation}
F_\chi
\equiv
K-\mathcal N_\chi[\gamma,\Sigma,K,\beta,\sigma,\pi_\sigma,Q,\chi,W^T,\psi,\pi_\psi],
\label{eq:Fchi}
\end{equation}
so that the longitudinal auxiliary equation is
\begin{equation}
\kappa_\theta \Delta_\gamma \chi = -F_\chi.
\label{eq:chieq}
\end{equation}
The old unweighted closure test failed because nonzero \(\mathfrak m_\chi=\int F_\chi\) produces the Coulomb asymptotic of the previous unit-lapse obstruction note.

\begin{definition}[Recorded unit-mass profile for the Coulomb split]
Fix once and for all a smooth compactly supported function
\begin{equation}
B_{\mathrm{ext}}\in C_c^\infty(\mathbb R^3)
\qquad\text{such that}\qquad
\int_{\mathbb R^3} B_{\mathrm{ext}}(x)\,dx = 1.
\label{eq:Bext}
\end{equation}
\end{definition}

\begin{definition}[Coulomb split of the longitudinal auxiliary variable]
On a slice with source monopole
\begin{equation}
\mathfrak m_\chi
\equiv
\int_{\mathbb R^3}F_\chi(x)\,dx,
\label{eq:mchi}
\end{equation}
define the \emph{Coulomb profile} \(\chi_{\mathrm{Coul}}\) and the \emph{regular part} \(\chi_{\mathrm{reg}}\) by
\begin{equation}
\chi=\chi_{\mathrm{Coul}}+\chi_{\mathrm{reg}},
\label{eq:split}
\end{equation}
where \(\chi_{\mathrm{Coul}}\) is the decaying solution of
\begin{equation}
\kappa_\theta \Delta \chi_{\mathrm{Coul}}
=
-\mathfrak m_\chi B_{\mathrm{ext}}
\label{eq:chiCoul}
\end{equation}
and \(\chi_{\mathrm{reg}}\) is the remainder.
\end{definition}

\begin{remark}
Definition 2 fixes only the low-frequency renormalization. It does not change the completion \(S_{\mathrm{min},4}\), the unit-lapse gauge, or the propagating Einstein--quartic-scalar variables.
\end{remark}

Since the present test concerns the old asymptotically flat Sobolev package, it is enough to examine the flat-background small-data model in which \(\Delta_\gamma\) is a perturbation of the Euclidean Laplacian. The decisive low-frequency issue is already visible in the Euclidean part.

\begin{proposition}[The split removes the explicit monopole from the flat model]
\label{prop:meanzero}
On the flat model \(\gamma=\delta\), the regular part satisfies
\begin{equation}
\kappa_\theta \Delta \chi_{\mathrm{reg}}
=
-F_{\chi,\mathrm{reg}},
\qquad
F_{\chi,\mathrm{reg}}
\equiv
F_\chi-\mathfrak m_\chi B_{\mathrm{ext}},
\label{eq:chireg}
\end{equation}
and the renormalized source obeys
\begin{equation}
\int_{\mathbb R^3} F_{\chi,\mathrm{reg}}(x)\,dx = 0.
\label{eq:meanzero}
\end{equation}
\end{proposition}

\begin{proof}
Subtract \eqref{eq:chiCoul} from the flat equation \(\kappa_\theta\Delta\chi=-F_\chi\). The mean-zero property \eqref{eq:meanzero} follows from \eqref{eq:Bext} and \eqref{eq:mchi}.
\end{proof}

\begin{remark}
Proposition \ref{prop:meanzero} is exactly what the Coulomb split is supposed to accomplish. It removes the explicit monopole from the elliptic source and therefore removes the \emph{recorded} \(1/|x|\) tail at the formal source level.
\end{remark}

\section{The Old Energy One Would Like to Reuse}

The old local/coercive architecture would now try to replay the same unweighted energy with \(\chi\) replaced by \(\chi_{\mathrm{reg}}\).

\begin{definition}[Coulomb-split unweighted test package]
Fix \(N\ge 6\). The \emph{Coulomb-split unweighted test package} is the direct old energy with the longitudinal variable replaced by its regular part:
\begin{align}
\mathfrak E_N^{\mathrm{ul,reg}}
\equiv\;&
\|\gamma-\delta\|_{H^N}^2
+ \|\Sigma\|_{H^{N-1}}^2
+ \|K\|_{H^{N-1}}^2
+ \|\sigma\|_{H^{N-1}}^2
+ \frac{1}{m_S^2}\|\pi_\sigma\|_{H^{N-1}}^2
\nonumber\\
&+
\frac{\lambda_4}{M_*^2}\|\Delta_\gamma \sigma\|_{H^{N-1}}^2
+ \mathcal E_N^{\mathrm{matter}}
+ \|\beta\|_{H^{N+1}}^2
+ \|Q\|_{H^N}^2
+ \|\chi_{\mathrm{reg}}\|_{H^N}^2
+ \|W^T\|_{H^N}^2.
\label{eq:Ereg}
\end{align}
\end{definition}

\begin{remark}
Definition 3 is the exact question the present note tests. If the Coulomb split repaired the old energy, then \(\|\chi_{\mathrm{reg}}\|_{H^N}\) would be controlled by the same kind of unweighted Sobolev norms that already appear in \eqref{eq:Ereg}, with no added weighted or homogeneous seminorm.
\end{remark}

\section{Low-Frequency Structure of the Regular Part}

The main issue is not the high-frequency elliptic estimate. It is the remaining low-frequency singularity of the inverse Laplacian.

\begin{proposition}[Fourier representation of the regular part]
\label{prop:fourier}
Assume the flat model of Proposition \ref{prop:meanzero}. Then
\begin{equation}
\widehat{\chi_{\mathrm{reg}}}(\xi)
=
-\frac{\widehat{F_{\chi,\mathrm{reg}}}(\xi)}{\kappa_\theta |\xi|^2},
\label{eq:fourier}
\end{equation}
and therefore
\begin{align}
\|\chi_{\mathrm{reg}}\|_{H^N}^2
\lesssim\;&
\int_{|\xi|\le 1}
|\xi|^{-4}\,|\widehat{F_{\chi,\mathrm{reg}}}(\xi)|^2\,d\xi
\nonumber\\
&+
\|F_{\chi,\mathrm{reg}}\|_{H^{N-2}}^2.
\label{eq:splitestimate}
\end{align}
Conversely, the low-frequency integral in \eqref{eq:splitestimate} is necessary for unweighted \(L^2\) control of \(\chi_{\mathrm{reg}}\).
\end{proposition}

\begin{proof}
Taking the Fourier transform of \eqref{eq:chireg} gives \eqref{eq:fourier}. Then
\[
\|\chi_{\mathrm{reg}}\|_{H^N}^2
\simeq
\int_{\mathbb R^3}
(1+|\xi|^2)^N |\widehat{\chi_{\mathrm{reg}}}(\xi)|^2\,d\xi.
\]
Insert \eqref{eq:fourier} and split the integral into \(|\xi|\le 1\) and \(|\xi|\ge 1\). On the high-frequency region,
\[
(1+|\xi|^2)^N |\xi|^{-4}
\lesssim
(1+|\xi|^2)^{N-2},
\]
which gives the \(H^{N-2}\) term. On the low-frequency region, the multiplier remains \(|\xi|^{-4}\), yielding the first term in \eqref{eq:splitestimate}. Since the \(H^N\)-norm contains the \(L^2\)-norm, no unweighted \(H^N\) estimate can avoid that low-frequency integral.
\end{proof}

\begin{definition}[Missing low-frequency seminorm]
For the Coulomb-split source define
\begin{equation}
\mathfrak L_\chi^2
\equiv
\int_{|\xi|\le 1}
|\xi|^{-4}\,|\widehat{F_{\chi,\mathrm{reg}}}(\xi)|^2\,d\xi.
\label{eq:Lchi}
\end{equation}
\end{definition}

\begin{remark}
The seminorm \(\mathfrak L_\chi\) is precisely what the old energy does \emph{not} contain. It is a low-frequency or homogeneous control quantity rather than an ordinary unweighted Sobolev norm.
\end{remark}

\section{Why Zero Mean Is Still Not Enough}

Subtracting the monopole forces \(\widehat{F_{\chi,\mathrm{reg}}}(0)=0\), but that does not control the full seminorm \eqref{eq:Lchi}.

\begin{proposition}[Zero mean does not restore the old elliptic estimate]
\label{prop:noestimate}
There is no constant \(C_N>0\) such that every smooth zero-mean source \(g\) satisfies
\begin{equation}
\|(-\Delta)^{-1}g\|_{H^N(\mathbb R^3)}
\le
C_N \|g\|_{H^{N-2}(\mathbb R^3)}.
\label{eq:falseestimate}
\end{equation}
In particular, zero mean alone does not close the slice estimate for \(\chi_{\mathrm{reg}}\) in the old unweighted energy.
\end{proposition}

\begin{proof}
Choose a nonzero \(\varphi\in C_c^\infty(\mathbb R^3)\) supported in the annulus
\[
\frac{1}{2}\le |\xi|\le 2.
\]
For \(j\ge 1\), define \(g_j\in \mathcal S(\mathbb R^3)\) by
\begin{equation}
\widehat{g_j}(\xi)
\equiv
2^{3j/2}\,\varphi(2^j\xi).
\label{eq:gj}
\end{equation}
Then \(\widehat{g_j}(0)=0\), so each \(g_j\) has zero mean. Since \(\widehat{g_j}\) is supported where \(|\xi|\simeq 2^{-j}\), one has
\[
\|g_j\|_{H^{N-2}}^2
\simeq
\int |\widehat{g_j}(\xi)|^2\,d\xi
\simeq
1,
\]
uniformly in \(j\). Let \(u_j=(-\Delta)^{-1}g_j\). Then
\[
\widehat{u_j}(\xi)
=
|\xi|^{-2}\widehat{g_j}(\xi),
\]
so on the support of \(\widehat{g_j}\),
\[
|\xi|^{-2}\simeq 2^{2j}.
\]
Therefore
\[
\|u_j\|_{L^2}^2
\simeq
2^{4j}\int |\widehat{g_j}(\xi)|^2\,d\xi
\simeq
2^{4j},
\]
which tends to infinity. Since the \(H^N\)-norm dominates the \(L^2\)-norm, \eqref{eq:falseestimate} cannot hold.
\end{proof}

\begin{corollary}[What the Coulomb split actually buys]
\label{cor:whatbuys}
The Coulomb split removes the explicit monopole contribution, but it does \emph{not} convert the longitudinal elliptic solve into an unweighted \(H^{N-2}\to H^N\) estimate. The remaining obstacle is the low-frequency seminorm \(\mathfrak L_\chi\), not the monopole itself.
\end{corollary}

\begin{proof}
Proposition \ref{prop:meanzero} gives the monopole subtraction. Proposition \ref{prop:noestimate} shows that the old unweighted elliptic estimate still fails.
\end{proof}

\section{Verdict for the Old Unweighted Package}

The previous proposition already decides the question the present note was meant to test.

\begin{theorem}[Failure of the Coulomb-split regular-part closure test]
\label{thm:failure}
Assume the fixed-completion unit-lapse spatial-harmonic formulation for \(S_{\mathrm{min},4}\) and apply the Coulomb split \eqref{eq:split} to the longitudinal auxiliary variable. Then, at the present recorded scope:
\begin{enumerate}
    \item the split removes the explicit monopole from the slice equation through Proposition \ref{prop:meanzero}
    \item nevertheless the regular part \(\chi_{\mathrm{reg}}\) does not satisfy the old unweighted slice estimate
    \[
    \|\chi_{\mathrm{reg}}\|_{H^N}
    \lesssim
    \|K\|_{H^{N-1}}+\cdots
    \]
    in terms of the unweighted Sobolev quantities already present in \eqref{eq:Ereg}
    \item therefore the Coulomb split does not presently repair the old unweighted local well-posedness, continuation, or coercive architecture for the \((K,\chi,\beta)\) package.
\end{enumerate}
\end{theorem}

\begin{proof}
Item (1) is Proposition \ref{prop:meanzero}. Item (2) follows from Proposition \ref{prop:noestimate}: even for smooth zero-mean sources, the inverse Poisson operator is not bounded from \(H^{N-2}\) to unweighted \(H^N\). Since the old energy \eqref{eq:Ereg} contains only unweighted Sobolev norms and no low-frequency seminorm such as \eqref{eq:Lchi}, it cannot close the elliptic estimate for \(\chi_{\mathrm{reg}}\). Item (3) then follows exactly as in the earlier obstruction note: if the slice-wise auxiliary estimate fails, the hyperbolic energy argument that depends on it cannot be replayed in the same function class.
\end{proof}

\begin{remark}
Theorem \ref{thm:failure} is a negative result for the \emph{old unweighted package}, not for Coulomb splitting in every possible formulation. If one enlarges the control norm by a weighted moment, a homogeneous low-frequency seminorm, or some equivalent replacement for \(\mathfrak L_\chi\), then a different theorem route may remain available.
\end{remark}

\section{What the Next Reformulation Must Record}

The old energy has now failed twice after the gauge-first reformulation: once because of the explicit monopole tail, and again because subtracting that tail still leaves an uncontrolled low-frequency inverse-Laplacian singularity.

\begin{corollary}[Next honest step after the Coulomb-split failure]
\label{cor:next}
The next honest fixed-completion reformulation is now the weighted or homogeneous function-space route in which the controlled longitudinal quantity is no longer \(\chi\) or \(\chi_{\mathrm{reg}}\) in plain unweighted \(H^N\), but rather a package that includes the missing low-frequency control represented by \(\mathfrak L_\chi\) or an equivalent weighted moment norm.
\end{corollary}

\begin{proof}
Theorem \ref{thm:failure} shows that the Coulomb split does not recover the old energy class. Proposition \ref{prop:fourier} identifies exactly what is missing: a low-frequency control quantity beyond unweighted Sobolev regularity. Therefore the next reformulation must change the function space rather than merely subtracting the monopole profile inside the old one.
\end{proof}

\section{Conclusion}

The propagated trace-monopole cancellation test did not close, so the next smallest repair was the Coulomb split. That test is now complete as well.

Its verdict is precise. Subtracting the explicit monopole profile does fix the first visible \(1/|x|\) tail at the slice equation level, but it does not restore the old unweighted closure package. The reason is structural: after monopole subtraction the regular part still solves a Poisson equation, and the inverse Poisson operator on \(\mathbb R^3\) retains a low-frequency \(|\xi|^{-2}\) singularity that is not controlled by zero mean plus the old unweighted \(H^{N-2}\) source norm.

Therefore the fixed-completion unit-lapse route has now exhausted the two smallest repairs compatible with the old unweighted energy. The next burden is the broader weighted or homogeneous function-space reformulation in which the longitudinal auxiliary sector is controlled by a norm that records the missing low-frequency information explicitly.

\end{document}
